\section{Discussion and open problems}

The exact contribution is an operator contract, not a replacement for thermodynamic extrapolation. Fixed-basis measurement leaks from the local permutation span; the subsequent independent gate twirl restores the post-gate bond space. Compression and Hilbert--Schmidt positivity organize the spectrum, and the $n=2$ Ising identification recovers an established benchmark~\cite{Bao2020}. General-$n$ ranks, exceptional sets and growth-rate regularity remain mathematical questions even where finite computations are accurate.

For phase identification, the revision replaces short-ring universality assignments by two complementary methods. Uniform MPS accesses infinite spatial extent but has finite bond dimension and optimizer systematics. Tilted Clifford SMC has positive weights and exact tableau updates but has finite-population, finite-time and genealogy errors. Neither method removes the need for a controlled extrapolation. At $n=3$ the completed $\chi=96$ points support a single transition near $.3049$; weak first order is favored but not established. At $n=4$ the four-state-Potts claim is withdrawn, not suspended pending another fit. At $n=5$ the raw gap ladder remains consistent with stronger first-order behavior, without a demonstrated positive asymptotic interface tension.

The $n!$ comparison organizes the sectors more accurately than $q=n$, but it is not a new theorem of universality. Ref.~\cite{Vasseur2019} already establishes permutation-spin and Potts-limit connections in its setting, and the regular bimodule is standard algebra. The contribution here is their disciplined use in this particular monitored-circuit model and the independent numerical tests they motivate. A real-moment changeover between the Ising rung and $n=3$ is conditional on the latter being first-order and on the existence and regularity of the interpolating phase boundary. No unique $n_*$ or replica CFT is derived. Fits toward $q=1$ remain conjectural and do not determine a quenched entropy, in accordance with Theorem~\ref{thm:no-go}.

The next computation is population-controlled tilted SMC, not merely a larger raw histogram. It should combine exact finite-time anchors near $.305$, population ladders and independent seeds, sufficient temporal buffers, blocked additive-density observables and a direct check that both candidate phases are sampled. Genealogical ESS and phase-switching histories are required alongside incremental ESS. Bias or non-mixing at coexistence may require tempering, trajectory-space moves or an improved guiding function. A larger MPS calculation should separately resolve two branches at the same $p$, quantify residuals and test correlation-length saturation; $\chi\ge200$ is a computational target, not a guarantee. The unchanged six-label MPO at $n=3,d=3$ is an economical next dimension test. A sector-aware $n=4$ MPS implementation and the 29-dimensional two-qubit Clifford fourth commutant are further tasks, not completed computations.

The analytic program also has explicit gates. A large-$d$ expansion must be derived for this monitored-circuit operator before an ordinary Potts result is imported. A Pirogov--Sinai or contour route, including the methods of Borgs and Imbrie~\cite{BorgsImbrie1989}, would require identified stable phases and verified contour/Peierls bounds; the presence of signed coordinates does not make the theorem automatically applicable. A cluster or loop representation exposing real $q$ as a parameter would be valuable, but is not constructed. The exact $n=2$ endpoint $p_c(d)\to1$ and a proposed $n\to\infty$ endpoint should not be conflated with a proof of noncommuting replica and dimension limits.

The record and complexity results retain independent standing. Appendix~\ref{app:complexity} concerns worst-case universal-circuit descriptions and promised conditional output or purity; it implies neither average-case Haar hardness nor hardness of a particular scalar record entropy. The exact replica interpolation identity remains valid at finite volume, but its limiting sufficient condition must control tilted variances extensively and uniformly, not just bound the variance of an intensive quantity. The no-freezing result requires an analytic nondegenerate branch and the stated cocycle hypotheses: neither all finite-state reducible processes nor an unverified full-Haar analyticity assertion follows merely from positivity. These qualifications are independent of the replica class dispute and are recorded rather than hidden by it.

For quenched Haar record statistics, the reported finite-size SCGFs are numerical objects with population, time, size and disorder uncertainty; they do not establish a thermodynamic freezing transition. The Clifford disorder tilt is a different object from its affine fixed-disorder trajectory SCGF. It can have nontrivial thermodynamic disorder statistics without contradicting that trajectory affinity. A derivation of a quenched continuum theory would have to control both replica limits and Born normalization, rather than infer it from the annealed ladder.

\section{Interpretive status of the record formalism}
\label{sec:interpretation}

The record can be read operationally as an ordered stream of outcomes. Its probabilities, entropy rates, LDPs and conditional states are specified within the circuit model. Calling the record an ``appearance-stream'' changes no result and establishes no ontology. Theorems about a stochastic process presupposing an event order do not derive time, consciousness, a universal observer, or the identity of a creator with creation.

The exact distinction is informational: symmetric unconditioned summaries need not determine Born-weighted conditional entanglement, while record-resolved information can. This does not derive the Born rule; that rule is a premise of the process. Nor does self-averaging mean that agents with different histories must converge to identical probabilities without common model and inference assumptions. It is a statement about specified random quantities and limits.

Interpretations that reproduce the same operational instrument reproduce the same record law. A modified-collapse dynamics, a different instrument, or a different probability rule need not do so, and is not covered by a theorem about the original law. There is no theorem here making Bohmian, Everettian, QBist or collapse descriptions equivalent in general. The phase-transition retraction therefore changes no interpretation-specific prediction; the old references to ``specific Potts classes'' and marginal corrections as established facts are removed.

The manuscript's commitments are checkability, explicit primitives, stable distinctions between theorem, numerical evidence and conjecture, and reproducible artifacts. Its former non-dual working assumption is not required or imported as a premise. Any separate philosophical account must be argued independently; neither a theological premise nor its negation is established by these circuit calculations. This section carries no evidential weight for the physics.

\section{Conclusion}
\label{sec:conclusion}

The paper establishes exact local leakage and post-gate closure, a compressed positive-semidefinite spectral description, rank certificates in their stated scope, and a recovery of the Bao--Choi--Altman two-replica Ising solution with closed spectrum. These are independent of the withdrawn numerical class assignments.

The corrected annealed phase diagram has a numerically located three-replica transition near $p_c^{(3)}=.3049(2)$, with simultaneous loss of sign and standard order in the sampled uniform-MPS scans through $\chi=96$. Weak first order is favored; continuity remains open. We retract the three- and four-state-Potts assignments and the argument from the dimension of one standard irrep. The regular bimodule and the conditional $n!$ ordered-state orbit supply the appropriate structural comparison without proving transition order. Four- and five-replica small-system spectra are retained as such; the $n=5$ first-order reading is plausible but not asymptotically certified. No converged coexistence length, nonzero latent heat, real-replica changeover, or $q\to1$ continuation is claimed.

Conditional-trajectory and conditional-purity decision remain PostBQP-complete under their promises. The record theory retains the exact LDP/multifractal and interpolation identities, numerical quenched Haar SCGFs, a correctly scoped no-freezing statement, and the Clifford dyadic closure through the third Haar moment. The new tilted stabilizer sampler independently reproduces finite-time anchors and reaches $L=128$ without signed weights or MPS truncation. Its large-size genealogies nevertheless collapse, so this pilot does not settle coexistence. The two-qubit fourth Clifford commutant has rank 29, not an independent stable-regime 30-label basis; its transfer operator and possible critical-point split are future computations.

The methodological result is not that exact spectra settle universality on small lattices, but that exact operators can be reused with independent, controlled limiting methods. The revised paper separates what is proved, what is reproduced, what is preferred, and what still requires a discriminating computation.
